% Complete finite restricted-decision, testing and product proofs.

\begin{annexitem}{annex:p17.randomized-decision-hulls}{thm:p17-randomized-decision-hulls}{Product mixtures of deterministic selections}
\begin{proof}
Let \(M\) be row stochastic with \(r\) rows and \(k\) columns. For each
map \(f:\{1,\ldots,r\}\to\{1,\ldots,k\}\), put
\[
                           w_f=\prod_{z=1}^r M_{z,f(z)}.
\]
The weights are nonnegative. Finite distributivity gives
\[
 \sum_f w_f=\prod_{z=1}^r\sum_{j=1}^kM_{zj}=1,\qquad
 \sum_{f:f(z)=j}w_f=M_{zj}\prod_{v\ne z}\sum_iM_{vi}=M_{zj}.
\]
An empty product in the second equality is one.
Consequently \(M=\sum_f w_fM^f\), so \(PM\) lies in the claimed convex
hull. Conversely, any convex combination of the \(M^f\) is row stochastic
and its product with \(P\) is the same convex combination of their products.
This proves the first identity, including \(k=1\).

For the second, take choices \(a_z\in\operatorname{conv}C\).
By definition, each is a convex combination of finitely many points of
\(C\). The union of those finite lists across the finitely many signals
is a finite list \(c_1,\ldots,c_N\). Add zero coefficients where needed
to write
\[
                   a_z=\sum_{j=1}^N\lambda_{zj}c_j,\qquad
                   \lambda_{zj}\geq0,\quad\sum_j\lambda_{zj}=1.
\]
For \(f:\{1,\ldots,r\}\to\{1,\ldots,N\}\), the deterministic choice
\(c_{f(z)}\) at signal \(z\) has a loss vector in \(B(P,C)\).
Give it weight \(\prod_z\lambda_{z,f(z)}\).
The preceding product calculation shows that the weighted loss vector
has state \(s\) coordinate
\[
             \sum_zp_{sz}\sum_j\lambda_{zj}(c_j)_s
                         =\sum_zp_{sz}(a_z)_s.
\]
Thus \(B(P,\operatorname{conv}C)\subseteq\operatorname{conv}B(P,C)\).
In the other direction, mix any finite collection of \(C\)-valued choices
with their convex weights at each signal. The resulting choice is in
\(\operatorname{conv}C\), and its loss vector is that convex mixture.
This proves equality without any compactness assumption on \(C\).

Finally, when \(a_1,\ldots,a_k\) generate \(\mathcal A\), a choice in
\(\mathcal A\) at each signal is represented by a row-stochastic matrix
\(M\) of their convex coefficients. Its risk vector is
\(\operatorname{diag}(PMD)\). Every such \(M\) gives allowable choices,
which proves both inclusions in the last formula.
\end{proof}
\end{annexitem}

\begin{annexitem}{annex:p17.finite-k-comparison}{thm:p17-finite-k-comparison}{Loss hulls, separating payoffs and posterior maxima}
\begin{proof}
\paragraph{From obtainable experiments to loss and payoff comparisons.}
Assume item 1. For a hull of \(k\) loss vectors with row matrix \(D\),
the preceding lemma expresses its risk set as
\(\{\operatorname{diag}(RD):R\in\mathcal P_{P,k}\}\).
Inclusion of the obtainable-experiment sets therefore gives item 2.
Repeated vectors handle a smaller nonempty generating list.

For any action set with at most \(k\) elements, take the convex hull
of its statewise negative payoff vectors. Every randomized policy selects
a point of that hull at each signal, and every hull point has such a
randomized representation. Item 2 therefore reproduces every statewise
payoff vector of a \(Q\)-policy by a \(P\)-policy. Multiplication by any
prior and summation preserve the expected payoff. This gives item 3;
using the specified full-support prior gives item 4.

\paragraph{A separating \(k\)-action problem gives the converse.}
The set \(\mathcal P_{P,k}\) is the continuous linear image of the
nonempty compact convex set of row-stochastic \(r\times k\) matrices.
It is therefore nonempty compact and convex, by the finite-dimensional
compactness proof in \cref{thm:p17-finite-separation-minimax}.
Suppose an \(n\times k\) matrix \(R\in\mathcal P_{Q,k}\) is absent from
\(\mathcal P_{P,k}\). The separation lemma supplies \(H\in\mathbb R^{n\times k}\)
with
\[
                 \langle H,R\rangle>
                     \max_{T\in\mathcal P_{P,k}}\langle H,T\rangle.
\]
Choose the \(k\) actions \(1,\ldots,k\) with payoffs
\(u(j,s)=H_{sj}/\pi_s\). Every denominator is positive.
The right side is exactly \(V(\pi,P,u)\): a policy on \(P\) is one
of the matrices \(M\) defining \(\mathcal P_{P,k}\).
Write \(R=QN\) with \(N\) row stochastic. Using \(N\) as the policy
on \(Q\) achieves \(\langle H,R\rangle\), contradicting item 4.
Hence all such \(R\) belong to \(\mathcal P_{P,k}\), proving item 1.

\paragraph{The source's loss-set route to the same converse.}
There is also a direct reconstruction of its item 2 to item 1 argument.
Fix \(R\in\mathcal P_{Q,k}\) and any nonempty closed bounded convex
loss set \(K\subseteq\mathbb R^n\).
A selection after \(R\) uses \(k\) points \(a_1,\ldots,a_k\in K\).
Let \(\mathcal A\) be their convex hull. Its resulting loss vector is
attainable after \(Q\), by first producing \(R\), and uses only choices
in \(\mathcal A\). Item 2 reproduces it after \(P\), still with choices
in \(\mathcal A\subseteq K\).
Thus \(B(R,K)\subseteq B(P,K)\) for every such \(K\).
The full finite comparison theorem gives \(P>R\), which means
\(R\in\mathcal P_{P,k}\). This verifies the source implication without
assuming that \(R\)'s \(k\) outputs already give a direct \(P\)-garbling.

\paragraph{Posterior maxima.}
For any list of at most \(k\) payoff vectors, define
\(\phi_u(p)=\max_j\sum_s p_su(j,s)\).
The previously proved value formula is
\[
                  V(\pi,P,u)=\int\phi_u\,d\mu_{P,\pi}.
\]
Every maximum of at most \(k\) linear functions on the simplex has
this form, because a linear function there is specified by its coefficients
at the states. Thus items 4 and 5 are equivalent.
The source's standard measure is the special case of the uniform prior;
the proof used only strict positivity of the chosen prior.

\paragraph{The source's finite-support posterior route.}
To recover its additional argument from item 5 to item 1, fix any
\(R\in\mathcal P_{Q,k}\). Its posterior law at \(\pi\) is supported
on at most \(k\) points, because \(R\) has \(k\) output labels.
Let \(\phi\) be the maximum of any finite nonempty list of linear
functions. At each positive-mass posterior of \(R\), select one
maximizing function from that list. There is at least one such posterior.
The maximum \(\psi\) of the selected functions has at most \(k\)
pieces, satisfies \(\psi\leq\phi\) everywhere, and agrees with \(\phi\)
at every posterior with positive \(R\)-mass. Therefore
\[
 \int\phi\,d\mu_{P,\pi}
 \geq\int\psi\,d\mu_{P,\pi}
 \geq\int\psi\,d\mu_{Q,\pi}
 \geq\int\psi\,d\mu_{R,\pi}
 =\int\phi\,d\mu_{R,\pi}.
\]
The middle comparisons use item 5 and the proved convex-order direction
for the garbling \(Q>R\), respectively. Thus every finite-action value
problem favors \(P\) over \(R\) at this full-support prior.
The full finite comparison theorem gives \(P>R\), proving item 1.

For the source's reverse posterior construction, assume item 1 and write
\(\phi=\max_{1\leq j\leq k}L_j\), with repetitions if fewer pieces occur.
At each positive-probability signal of \(Q\), choose the least index
maximizing \(L_j\) at that signal's posterior. Assign a zero-probability
signal arbitrarily. This deterministic partition produces
\(R\in\mathcal P^*_{Q,k}\).
At each positive-mass output \(j\) of \(R\), its posterior is the
conditional average of the \(Q\)-posteriors in that cell.
Each difference \(L_j-L_i\) is nonnegative at all of them, so it is
nonnegative at their average. Hence \(L_j\) is still a maximizer at
the \(R\)-posterior. Linearity and regrouping the cell probabilities give
\[
                       \int\phi\,d\mu_{Q,\pi}
                              =\int\phi\,d\mu_{R,\pi}.
\]
Null cells contribute zero. Item 1 gives \(R\in\mathcal P_{P,k}\),
so \(P>R\) and convex-order monotonicity gives item 5.
This supplies the source's partition and supporting-function arguments
with every finite support and tie case stated.
\end{proof}
\end{annexitem}

\begin{annexitem}{annex:p17.k-comparison-hierarchy}{thm:p17-k-comparison-hierarchy}{The restricted preorder and a sufficient output bound}
\begin{proof}
Reflexivity and transitivity follow from reflexivity and transitivity
of set inclusion for \(\mathcal P_{P,k}\). For \(k=1\), the only
row-stochastic processing matrix has every entry one, and every resulting
experiment is the same \(n\times1\) column of ones. This proves the
universal \(>_1\) equivalence.

If \(PM=Q\), every \(k\)-output transformation \(QN\) equals \(P(MN)\),
with \(MN\) row stochastic. Thus full dominance implies every restricted
dominance. Any payoff problem with at most \(k\) actions is also a problem
with at most \(k+1\) actions, so the equivalent payoff characterization
gives \(>_{k+1}\) implies \(>_k\).

Suppose \(Q\) has \(t\leq k\) outputs. Embed its outputs into the first
\(t\) of the \(k\) labels. The resulting \(Q'\) consists of \(Q\) and
\(k-t\) zero columns and belongs to \(\mathcal P_{Q,k}\).
If \(P>_kQ\), some matrix produces \(Q'\) from \(P\).
Map the first \(t\) output labels back to themselves and map each
remaining label to label 1. This is a row-stochastic deterministic
postprocessor and sends \(Q'\) to \(Q\), since the added columns are zero.
Composition gives \(P>Q\). The reverse implication was already proved.
Taking \(k=t\) also proves the all-finite-bounds assertion.
\end{proof}
\end{annexitem}

\begin{annexitem}{annex:p17.convex-hinge-decomposition}{thm:p17-convex-hinge-decomposition}{All breakpoints and their slope increments}
\begin{proof}
Write the restricted linear functions as affine functions
\(a_i+b_ix\). Remove identical duplicates.
Two different affine functions with different slopes intersect at at most
one point; two with equal slopes never intersect.
List in increasing order all distinct intersection points strictly
inside \([0,1]\), and append the endpoints 0 and 1.
Between two consecutive points the ordering of every pair of affine
functions is constant: a sign change of their affine difference would
force a zero there. Hence their maximum is one fixed affine function on
each such open interval. The maximum is continuous, being the maximum
of finitely many continuous functions, and so agrees with the corresponding
affine limits at interval endpoints.

The maximum is convex: each affine value at a convex combination equals
the corresponding combination of its values and is bounded above by the
combination of the two maxima.
If \(t\) is an interior partition point with adjacent slopes \(b_-\)
and \(b_+\), choose a sufficiently small \(h>0\) lying within its two
adjacent intervals. Convexity at their midpoint yields
\[
  2\phi(t)\leq\phi(t-h)+\phi(t+h)
             =2\phi(t)+h(b_+-b_-).
\]
Therefore \(b_+\geq b_-\).
Keep all interior partition points \(t_1,\ldots,t_m\), allowing a
zero slope increment at a point where the maximizing affine expression
does not change. Let \(b\) be the first interval's slope, \(a=\phi(0)\),
and let \(d_j\) be the right slope minus the left slope at \(t_j\).
These coefficients are nonnegative.

The proposed expression has value \(a\) at zero, slope \(b\) on the
first interval, and gains slope \(d_j\) at each \(t_j\).
It therefore has exactly the slope of \(\phi\) on every interval.
Both functions are continuous. Starting with equality on the first
interval and propagating across the finitely many shared endpoints
proves equality everywhere, including 0 and 1.
If there are no interior intersections the same argument gives the
affine formula with an empty sum.

On \(p_1=x,p_2=1-x\), the expression
\((1-t)p_1-tp_2\) equals \(x-t\).
Also \(a+bx=(a+b)p_1+ap_2\).
This proves both simplex linearity assertions and completes the decomposition.
\end{proof}
\end{annexitem}

\begin{annexitem}{annex:p17.binary-two-decisions}{thm:p17-binary-two-decisions}{From two linear pieces to every finite decision problem}
\begin{proof}
Full dominance implies the two-decision comparison by the hierarchy.
For the converse choose the uniform prior, or any other fixed full-support
prior. The two posterior laws have the same barycenter by
\cref{thm:p17-posterior-splitting}. In particular their means of
\(x=p_1\) agree and their masses are one.

Assume \(P>_2Q\). The restricted comparison theorem orders the expectations
of every hinge \((x-t)_+\), because it is the maximum of two linear
functions on the simplex. Their expectations of every affine expression
\(a+bx\) are equal by the common mean and total mass.
For any finite-action problem, its posterior value function is a maximum
of finitely many affine functions of \(x\). Apply the preceding
decomposition. Multiplying the hinge inequalities by their nonnegative
coefficients, summing, and adding the common affine expectation gives
the expected-value inequality for that problem.
Thus every finite payoff problem is dominated at the fixed full-support
prior. The full finite Blackwell theorem gives \(P>Q\).
This proof needs no approximation by infinitely many affine pieces.
\end{proof}
\end{annexitem}

\begin{annexitem}{annex:p17.testing-error-region}{thm:p17-testing-error-region}{All vertical sections and the comparison by their lower boundary}
\begin{proof}
\paragraph{The region and its sections.}
The error map is affine and continuous on \([0,1]^r\).
This cube is the convex hull of its finitely many zero-one vertices,
as follows from the same product weights used for deterministic choices.
Its image is therefore the convex hull of finitely many points: a compact
convex polygon, permitting the degenerate line-segment case.
The all-zero and all-one tests give \((0,1)\) and \((1,0)\).
Complementing every coordinate of a test sends its errors to
\((1-a,1-b)\), which proves central symmetry.

For every \(t\in[0,1]\), the constant test \(\delta_z=t\) has errors
\((t,1-t)\). The intersection of the region with the vertical line at
\(t\) is therefore nonempty, compact and convex, and hence a closed
interval of second coordinates. Its minimum is \(f_P(t)\).
Central symmetry sends this interval to the one at \(1-t\) and reverses
its vertical coordinate, so its maximum is \(1-f_P(1-t)\).
This proves the exact full-region formula.

For sizes \(t_0,t_1\), mix tests attaining their minima with weights
\(1-\lambda,\lambda\). The mixed test has first error
\((1-\lambda)t_0+\lambda t_1\) and second error
\((1-\lambda)f_P(t_0)+\lambda f_P(t_1)\).
Minimality proves the convexity inequality for \(f_P\).

For \(0\leq t<t'\leq1\), mix an optimal size-\(t\) test with the
always-reject test using weight \(\lambda=(t'-t)/(1-t)\) on the latter.
This is defined because \(t<1\), and its errors are
\((t',(1-\lambda)f_P(t))\).
Thus \(f_P(t')\leq(1-\lambda)f_P(t)\leq f_P(t)\).
The always-reject test also proves \(f_P(1)=0\).

An exact-size-\(t\) optimal test is feasible under a size-at-most-\(t\)
constraint. Conversely any test of size \(a\leq t\) has second error
at least \(f_P(a)\geq f_P(t)\). This proves equality of the two optima
and their attainment.

\paragraph{The region is a source risk set.}
Let \(A\) be the line segment joining \((0,1)\) and \((1,0)\).
A choice on this segment at signal \(z\) is exactly
\((\delta_z,1-\delta_z)\). Its statewise risk is
\((a_P(\delta),b_P(\delta))\), so \(\mathcal R(P)=B(P,A)\).

If \(P>Q\), all risk-set comparisons hold, including
\(\mathcal R(P)\supseteq\mathcal R(Q)\).
Conversely assume this inclusion. Any segment \(A'\) in \(\mathbb R^2\),
with endpoints \(c,d\), is the image of \(A\) under the coordinatewise
affine map
\[
 L(v_1,v_2)=
   \bigl(c_1+(d_1-c_1)v_1,\ d_2+(c_2-d_2)v_2\bigr).
\]
At \((v_1,v_2)=(\delta,1-\delta)\), this equals
\(c+\delta(d-c)\). The map is allowed to be singular; coincident
endpoints give a single point. Since each experiment row sums to one,
applying \(L\) to a risk vector gives the risk vector of the transformed
choices. Thus \(LB(P,A)=B(P,A')\), and the same holds for \(Q\).
The assumed inclusion therefore gives the comparison for every hull
of at most two loss points. This is \(P>_2Q\), hence \(P>Q\) by the
two-state theorem.

Finally, inclusion of regions implies \(f_P(t)\leq f_Q(t)\) by taking
the point minimizing the \(Q\)-section. If those inequalities hold at
all \(t\), they also give
\(1-f_P(1-t)\geq1-f_Q(1-t)\). The full-region formula therefore puts
every \(Q\)-section inside the corresponding \(P\)-section.
This proves the entire chain of equivalences.
\end{proof}
\end{annexitem}

\begin{annexitem}{annex:p17.finite-likelihood-optimality}{thm:p17-finite-likelihood-optimality}{The threshold exchange inequality with all zero masses}
\begin{proof}
Rejecting at a signal with \(p_z=0<q_z\) increases power without affecting
size. It is therefore optimal to reject there with probability one.
A signal with both probabilities zero contributes nothing.
The remaining signals have positive \(p_z\), and their \(p_z\) sum is one.
Consequently their cumulative sums increase strictly from zero to one,
so the selected first index \(j\) exists for \(t<1\).
Its preceding sum satisfies \(A_{j-1}\leq t<A_j\); hence the proposed
fraction lies in \([0,1)\). The test \(\delta^*\) has size exactly \(t\)
and power \(q_*+B_{j-1}+\ell_{z_j}(t-A_{j-1})\).

Put \(c=\ell_{z_j}\geq0\). For a preceding positive-\(p\) signal,
\(q_z-cp_z\geq0\) and \(\delta_z^*=1\).
For a later one, \(q_z-cp_z\leq0\) and \(\delta_z^*=0\).
At the selected signal the coefficient is zero, regardless of its
randomization. At a zero-\(p\), positive-\(q\) signal the coefficient
is positive and the test rejects. At a double-zero signal the coefficient
is zero. Thus for any competing test \(\delta\),
\[
                        \sum_z(q_z-cp_z)(\delta_z-\delta_z^*)\leq0.
\]
If its size is at most \(t\), then
\[
 \sum_zq_z(\delta_z-\delta_z^*)
   \leq c\sum_zp_z(\delta_z-\delta_z^*)\leq0.
\]
The proposed test maximizes power and therefore minimizes the second
error even under the size-at-most constraint. This proves the formula
for \(t<1\). At \(t=1\), rejection everywhere gives second error zero,
which no test can improve.

In a tied ratio block, replacing one allocation of the available
first-kind mass by another changes power by the common ratio times the
change in allocated mass. That change is zero at a fixed total size.
Thus every tie ordering gives the same curve.

At size zero, nonnegativity forces \(\delta_z=0\) at every positive-\(p\)
signal. The best possible power is exactly \(q_*\), achieved by rejecting
at all null-impossible alternative signals. Hence \(f_P(0)=1-q_*\).
This also explains why an endpoint of the full region need not lie
on its lower boundary.
\end{proof}
\end{annexitem}

\begin{annexitem}{annex:p17.repeated-experiment-comparison}{thm:p17-repeated-experiment-comparison}{A finite product certificate for each sample size}
\begin{proof}
Every entry of \(P^{[m]}\) is nonnegative, and for a fixed state
\[
                    \sum_{z_1,\ldots,z_m}\prod_i p_{s,z_i}
                            =\prod_i\sum_{z_i}p_{s,z_i}=1.
\]
Thus it is an experiment. The same calculation at every fixed input
tuple \(z\) shows that \(M^{[m]}\) is row stochastic.
For each state and output tuple, finite distributivity gives
\[
 \begin{aligned}
 (P^{[m]}M^{[m]})_{s,y}
   &=\sum_{z_1,\ldots,z_m}\prod_{i=1}^m
                              p_{s,z_i}M_{z_i,y_i}\\
   &=\prod_{i=1}^m\sum_{z_i}p_{s,z_i}M_{z_i,y_i}\\
   &=\prod_{i=1}^m q_{s,y_i}=Q^{[m]}_{s,y}.
 \end{aligned}
\]
This is the claimed garbling certificate. The state is held fixed in
each product; averaging different states separately at each observation
would describe a different experiment.

Statewise policy simulation now proves all decision-value comparisons
for these repeated experiments. In the binary case the assumed
one-observation curve comparison gives \(P>Q\) by
\cref{thm:p17-testing-error-region}. Apply the product certificate,
then the same theorem to \(P^{[m]},Q^{[m]}\).
This proves the error-curve inequality for every positive sample size.
\end{proof}
\end{annexitem}

\begin{annexitem}{annex:p17.testing-calculations-boundaries}{thm:p17-testing-calculations-boundaries}{Explicit curves and a failure without conditional independence}
\begin{proof}
For \(P\), the ratios \(q_z/p_z\) are \(1/4\) at signal 1 and
\(4\) at signal 2. Reject signal 2 first. Its null mass is \(1/5\)
and alternative mass \(4/5\); thus power is \(4t\) up to size \(1/5\).
Past that size, power is \(4/5+(t-1/5)/4\), giving precisely the two
pieces stated for \(f_P\).
For \(Q\), the corresponding first selected signal has null mass \(7/20\),
alternative mass \(13/20\) and ratio \(13/7\).
The other ratio is \(7/13\). The same construction gives the displayed
\(f_Q\). At \(t=1/5\), its first piece is
\(1-13/35=22/35\), whereas \(f_P(1/5)=1/5\).

The independent two-observation masses are
\[
 \begin{array}{c|cccc}
 (z_1,z_2)&(2,2)&(1,2)&(2,1)&(1,1)\\\hline
 \text{state 1}&1/25&4/25&4/25&16/25\\
 \text{state 2}&16/25&4/25&4/25&1/25\\
 q_z/p_z&16&1&1&1/16 .
 \end{array}
\]
The first ratio block gives power \(16t\) until \(t=1/25\).
The tied middle block then gives power
\(16/25+(t-1/25)=3/5+t\) until cumulative null mass \(9/25\).
At that size, power is \(24/25\). The last block adds
\((t-9/25)/16\), whose complementary error is \((1-t)/16\).
These calculations give all three pieces and agree at their boundaries.
At \(t=1/10\), the middle piece has error
\(2/5-1/10=3/10\), and the one-observation first piece has error \(3/5\).

For a single signal, a deterministic test either always rejects or
always accepts. Its size is respectively 1 or 0; size \(1/2\), for
example, is unavailable. Randomizing at the signal with probability
\(t\) gives size exactly \(t\) and error \(1-t\).
With null law concentrated on signal 1 and alternative law concentrated
on signal 2, reject signal 2 and randomize with probability \(t\) at
signal 1. This has size \(t\) and second error zero for every \(t\).
It also verifies the singular size-zero boundary.

For the dependence counterexample, let both signals be bits. Define the
joint experiments \(A,B\) by the following conditional tables:
\[
 \begin{array}{c|cccc}
 (z_1,z_2)&(0,0)&(0,1)&(1,0)&(1,1)\\\hline
 A,\ \text{state 1}&1/2&0&0&1/2\\
 A,\ \text{state 2}&1/2&0&0&1/2\\
 B,\ \text{state 1}&1/2&0&0&1/2\\
 B,\ \text{state 2}&0&1/2&1/2&0 .
 \end{array}
\]
Every individual coordinate in every row is a fair bit. Thus the
one-observation marginal experiments are identical in both states
and are uninformative. But the rows of the joint experiment \(A\)
are identical, whereas those of \(B\) have disjoint supports.
Under equal state probabilities and unit payoff for a correct state guess,
\(A\) has optimal value \(1/2\), while \(B\) has value one: use equality
of its two bits to identify state 1 and inequality to identify state 2.
A garbling from \(A\) would preserve equality of its state rows, so
cannot produce \(B\). This proves that the equal marginal comparison
does not control an unspecified dependence structure.
\end{proof}
\end{annexitem}
